% ============================================================================
% SEC04.TEX - Section 11.4: Behavior Patterns
% ============================================================================

\section{Behavior Patterns}

\begin{learningoutcomes}
\item Mengimplementasikan patrol behavior
\item Mengimplementasikan chase behavior
\item Mengimplementasikan attack behavior
\end{learningoutcomes}

\subsection{Patrol Behavior}

\begin{lstlisting}[language={[Sharp]C}]
public class PatrolBehavior : MonoBehaviour
{
    public Transform[] waypoints;
    public float waitTime = 2f;
    
    private int currentWaypoint = 0;
    private NavMeshAgent agent;
    private bool waiting = false;
    
    void Start()
    {
        agent = GetComponent<NavMeshAgent>();
        MoveToWaypoint();
    }
    
    void Update()
    {
        if (!waiting && agent.remainingDistance < 0.5f)
        {
            StartCoroutine(WaitAtWaypoint());
        }
    }
    
    void MoveToWaypoint()
    {
        agent.SetDestination(waypoints[currentWaypoint].position);
    }
    
    IEnumerator WaitAtWaypoint()
    {
        waiting = true;
        yield return new WaitForSeconds(waitTime);
        currentWaypoint = (currentWaypoint + 1) % waypoints.Length;
        MoveToWaypoint();
        waiting = false;
    }
}
\end{lstlisting}

\subsection{Chase Behavior}

\begin{lstlisting}[language={[Sharp]C}]
public class ChaseBehavior : MonoBehaviour
{
    public Transform target;
    public float chaseSpeed = 5f;
    private NavMeshAgent agent;
    
    void Start()
    {
        agent = GetComponent<NavMeshAgent>();
        agent.speed = chaseSpeed;
    }
    
    void Update()
    {
        if (target != null)
        {
            agent.SetDestination(target.position);
        }
    }
}
\end{lstlisting}

\begin{catatan}
Behavior patterns dapat dikombinasikan dengan state machine untuk AI yang kompleks. Gunakan coroutine untuk timing behavior seperti wait atau delay.
\end{catatan}
